% !TEX encoding = UTF-8 Unicode
% !TEX TS-program = pdflatex

\documentclass[%
    12pt, % font size
    oneside, % if it should be done for printing on two-sides or not
    a4paper, % format of paper
    notitlepage, % to remove the title page
]{article}

%%%%%%%%%%%%%%%%%%%%%%%%%%%%%%%%%%%%%%%%%%%%%%%%%%%%
\usepackage[english]{babel}
\usepackage[utf8]{inputenc}
\usepackage[T1]{fontenc}
\usepackage{lmodern}
\usepackage{hyperref} % must be loaded before glossaries-extra
\usepackage{setspace}
% bibliography
\usepackage[hyperref=true,backref=true,backend=biber,maxbibnames=9,maxcitenames=2,style=numeric,citestyle=numeric,sorting=none]{biblatex} % hyperref uses links, backref goes back to citations, uses biber as backend, with 9 names at most in bibliography and 2 in citations, citing using numbers, and sorting in citation order
% sorting can be also ydnt for year descending, name, title or ynt for ascending year

\usepackage{adjustbox} % to resize boxes by keeping the same aspect ratio
\usepackage{algorithm} % algorithm environment
\usepackage{algpseudocode} % improved pseudo-code
\usepackage{amsfonts}               %  AMS mathematical fonts
\usepackage{amsmath}
\usepackage{amssymb}                %  AMS mathematical symbols
\usepackage{bm}                     %  black/bold mathematical symbols
\usepackage{booktabs}               %  better tables
\usepackage[labelfont=bf]{caption} % font=footnotesize % to have reduced caption font size
\usepackage{csquotes}
\usepackage{enumitem} %left align the bulleted points

\usepackage{fancyhdr}

\usepackage[a4paper, left=2.54cm, right=2.54cm, top=2.54cm, bottom=2.54cm, bindingoffset=1.9cm]{geometry}
\usepackage{glossaries} % to use acronyms and glossary, it has also glossaries-extra as extension, but commands are different
\usepackage[%
    toc, % puts the link in the ToC
    %record, % to use bib2gls
    abbreviations, % to load abbreviations / acronyms
    nonumberlist, % to avoid printing the numbers of the references in the acronyms page
]{glossaries-extra}

\usepackage{graphicx}               %  post-script images
%\usepackage{iwona} % extra fonts, substitute standard ones
\usepackage{listings} % to insert formatted code
\usepackage{lipsum} % for lorem ipsum text, not needed in the real work
\usepackage{makecell} % to change dimensions of cells, for math cases
\usepackage{mathtools} % for additional commands
\usepackage{mfirstuc} % to have capitalization capabilities
\usepackage[final]{microtype}      % microtypography, final lets latex use it also in bibliography
\usepackage{multirow} % to allow for cells covering more than 1 row in tables
\usepackage{nicefrac}       % compact symbols for 1/2, etc.
%\usepackage[lofdepth,lotdepth]{subfig}
\usepackage{ragged2e} % for justifying text
\usepackage{siunitx} % support for SI units of measurement and number typesetting
\usepackage{subfig}
\usepackage{svg} % for svg support, works only if inkscape is installed, default for Overleaf v2
%\usepackage{subfigure}              %  subfigure compatibility, can be removed if subfig
\usepackage{tabularx} % equal-width columns in tables
\usepackage{textcomp} % extra fonts and symbols
\usepackage{url}            % simple URL typesetting
\usepackage{verbatim} % for extended verbatim support
\usepackage{xcolor} % to define colors and use standard CSS names add dvipsnames as option, but it clashes with xcolor loaded in toptesi, pay attention that if it goes in conflict with tikz/beamer, simply use \documentclass[usenames,dvipsnames]{beamer}, along with other custom options when defining the document class

\usepackage{titlesec}
\usepackage{longtable}
\usepackage{subcaption}
\usepackage{threeparttable}

\usepackage{calc} % needed for complex computationsin latex language
% \usepackage[a4paper]{geometry} % needed to change margins


% configuration for glossaries
% convert and load converted glossaries in .tex ,format from .bib
\setabbreviationstyle{long-short-desc} % style before loading resources
% this command sets the style to title for long names of acronyms only in the glossary description, leading to capitalized first-letter for all words
% \glssetcategoryattribute{\glsxtrabbrvtype}{glossname}{capitalisewords} % doesn't work
% resources to load if using a bib file with bib2gls
%\GlsXtrLoadResources[%
% src={glossaries}, % name of the file without extension
% selection=all, % select all the entries
%]
% not needed
%\newglossary*{abbreviation}{Acronyms} % to change the name of this glossary for acronyms

%\renewcommand{\glsclearpage}{\paginavuota} % to allow glossaries to clear pages, done manually is better

% setup for hyperref
\hypersetup{%
    pdfpagemode={UseOutlines},
    bookmarksopen,
    pdfstartview={FitH},
    colorlinks,
    linkcolor={black}, % it is suggested to keep them black, since when printing it it costs per page, and if they have color it's twice the price per page
    citecolor={black},
    urlcolor={black}
  }
%

% setup for svg
\svgsetup{%
    inkscapeformat=pdf, % to force usage of PDF
    inkscapelatex=false, % to disable latex rendering of text, produces errors
}

% setup for siunitx, it does not work in the summary
\sisetup{%
    detect-all, % to use the same font as for writing when using \num
    mode=text, % to allow it to work also in math mode
    group-separator = {,}, % separator for number grouping
    group-minimum-digits = 3, % minimum number of digits a number must have to be grouped in 3-digit groups
}

% listings colours
\definecolor{rulecolor}{rgb}{0,0,0}
\definecolor{commentcolor}{rgb}{0,0.6,0}
\definecolor{linenumbercolor}{rgb}{0.5,0.5,0.5}
\definecolor{keywordcolor}{rgb}{0,0,0.95}
\definecolor{backcolor}{rgb}{1,1,1}%{0.95,0.95,0.92}
\definecolor{stringcolor}{rgb}{0.58,0,0.82}

% setup for lstlisting
\lstset{ %
	backgroundcolor=\color{backcolor},   % choose the background color; you must add \usepackage{color} or \usepackage{xcolor}; should come as last argument
	basicstyle=\footnotesize,        % the size of the fonts that are used for the code
	breakatwhitespace=false,         % sets if automatic breaks should only happen at whitespace
	breaklines=true,                 % sets automatic line breaking
	captionpos=t,                    % sets the caption-position to bottom
	commentstyle=\color{commentcolor},    % comment style
	extendedchars=true,              % lets you use non-ASCII characters; for 8-bits encodings only, does not work with UTF-8
	frame=single,	                   % adds a frame around the code
	keepspaces=true,                 % keeps spaces in text, useful for keeping indentation of code (possibly needs columns=flexible)
	keywordstyle=\color{keywordcolor},       % keyword style
	%language=VHDL,                 % the language of the code
	numbers=left,                    % where to put the line-numbers; possible values are (none, left, right)
	numbersep=5pt,                   % how far the line-numbers are from the code
	numberstyle=\tiny\color{linenumbercolor}, % the style that is used for the line-numbers
	rulecolor=\color{rulecolor},         % if not set, the frame-color may be changed on line-breaks within not-black text (e.g. comments (green here))
	showspaces=false,                % show spaces everywhere adding particular underscores; it overrides 'showstringspaces'
	showstringspaces=false,          % underline spaces within strings only
	showtabs=false,                  % show tabs within strings adding particular underscores
	stepnumber=1,                    % the step between two line-numbers. If it's 1, each line will be numbered
	stringstyle=\color{stringcolor},     % string literal style
	tabsize=4,	                   % sets default tabsize to 2 spaces
	title=\lstname,                   % show the filename of files included with \lstinputlisting; also try caption instead of title
	inputencoding=utf8,
	literate=
	{á}{{\'a}}1 {é}{{\'e}}1 {í}{{\'i}}1 {ó}{{\'o}}1 {ú}{{\'u}}1
	{Á}{{\'A}}1 {É}{{\'E}}1 {Í}{{\'I}}1 {Ó}{{\'O}}1 {Ú}{{\'U}}1
	{à}{{\`a}}1 {è}{{\`e}}1 {ì}{{\`i}}1 {ò}{{\`o}}1 {ù}{{\`u}}1
	{À}{{\`A}}1 {È}{{\'E}}1 {Ì}{{\`I}}1 {Ò}{{\`O}}1 {Ù}{{\`U}}1
	{ä}{{\"a}}1 {ë}{{\"e}}1 {ï}{{\"i}}1 {ö}{{\"o}}1 {ü}{{\"u}}1
	{Ä}{{\"A}}1 {Ë}{{\"E}}1 {Ï}{{\"I}}1 {Ö}{{\"O}}1 {Ü}{{\"U}}1
	{â}{{\^a}}1 {ê}{{\^e}}1 {î}{{\^i}}1 {ô}{{\^o}}1 {û}{{\^u}}1
	{Â}{{\^A}}1 {Ê}{{\^E}}1 {Î}{{\^I}}1 {Ô}{{\^O}}1 {Û}{{\^U}}1
	{œ}{{\oe}}1 {Œ}{{\OE}}1 {æ}{{\ae}}1 {Æ}{{\AE}}1 {ß}{{\ss}}1
	{ű}{{\H{u}}}1 {Ű}{{\H{U}}}1 {ő}{{\H{o}}}1 {Ő}{{\H{O}}}1
	{ç}{{\c c}}1 {Ç}{{\c C}}1 {ø}{{\o}}1 {å}{{\r a}}1 {Å}{{\r A}}1
	{€}{{\euro}}1 {£}{{\pounds}}1 {«}{{\guillemotleft}}1
	{»}{{\guillemotright}}1 {ñ}{{\~n}}1 {Ñ}{{\~N}}1 {¿}{{?`}}1
}


% biblatex setup
% generally 9000 is ok, if higher than 10000 it's bad
% If you want to break on URL numbers
\setcounter{biburlnumpenalty}{9000}
% If you want to break on URL lower case letters
\setcounter{biburllcpenalty}{9000}
% If you want to break on URL UPPER CASE letters
\setcounter{biburlucpenalty}{9000}

\definecolor{mGreen}{rgb}{0,0.6,0}
\definecolor{mGray}{rgb}{0.5,0.5,0.5}
\definecolor{mPurple}{rgb}{0.58,0,0.82}
\definecolor{backgroundColour}{rgb}{0.95,0.95,0.92}

\lstdefinestyle{CStyle}{
    backgroundcolor=\color{backgroundColour},   
    commentstyle=\color{mGreen},
    keywordstyle=\color{magenta},
    numberstyle=\tiny\color{mGray},
    stringstyle=\color{mPurple},
    basicstyle=\footnotesize\ttfamily,
    breakatwhitespace=false,         
    breaklines=true,                 
    captionpos=b,                    
    keepspaces=true,                 
    numbers=left,                    
    numbersep=5pt,                  
    showspaces=false,                
    showstringspaces=false,
    showtabs=false,                  
    tabsize=2,
    language=C,
    lineskip=-0.1ex
}

% \setlength\textwidth{7in}
% \setlength\textheight{10in}
% \setlength\oddsidemargin{(\paperwidth-\textwidth)/2 - 1in}
% \setlength\topmargin{(\paperheight-\textheight
% -\headheight-\headsep-\footskip)/2 - 1in}

% this configuration is close to TopTesi in width
% \newgeometry{a4paper,top=3cm,bottom=3cm,left=3cm,right=3cm,%
% heightrounded}
% margins as for libreoffice writer
\newgeometry{top=2cm,bottom=2cm,left=2cm,right=2cm,%
heightrounded}


\lstdefinestyle{CStyle}{
    backgroundcolor=\color{backgroundColour},   
    commentstyle=\color{mGreen},
    keywordstyle=\color{magenta},
    numberstyle=\tiny\color{mGray},
    stringstyle=\color{mPurple},
    basicstyle=\footnotesize\ttfamily,
    breakatwhitespace=false,         
    breaklines=true,                 
    captionpos=b,                    
    keepspaces=true,                 
    numbers=left,                    
    numbersep=5pt,                  
    showspaces=false,                
    showstringspaces=false,
    showtabs=false,                  
    tabsize=2,
    language=C,
    lineskip=-0.1ex
}



% how to change Contents to Table of Contents
\addto\captionsenglish{% Replace "english" with the language you use
  \renewcommand{\contentsname}%
    {Table of Contents}%
}

% to change the name of Abbreviations to Acronyms
% not needed if use use entry types and define those
% \renewcommand{\abbreviationsname}{Acronyms}

% to allow line comments in algorithms
\algnewcommand{\LineComment}[1]{\State \(\triangleright\) #1}

% to declare abs and norm
\DeclarePairedDelimiter\abs{\lvert}{\rvert}%
\DeclarePairedDelimiter\norm{\lVert}{\rVert}%

% Swap the definition of \abs* and \norm*, so that \abs
% and \norm resizes the size of the brackets, and the 
% starred version does not.
\makeatletter
\let\oldabs\abs
\def\abs{\@ifstar{\oldabs}{\oldabs*}}
%
\let\oldnorm\norm
\def\norm{\@ifstar{\oldnorm}{\oldnorm*}}
\makeatother


\lstdefinestyle{CStyle}{
    backgroundcolor=\color{backgroundColour},   
    commentstyle=\color{mGreen},
    keywordstyle=\color{magenta},
    numberstyle=\tiny\color{mGray},
    stringstyle=\color{mPurple},
    basicstyle=\scriptsize\ttfamily,
    breakatwhitespace=false,         
    breaklines=true,                 
    captionpos=b,                    
    keepspaces=true,                 
    numbers=left,                    
    numbersep=5pt,                  
    showspaces=false,                
    showstringspaces=false,
    showtabs=false,                  
    tabsize=2,
    language=C,
    lineskip=-0.1ex
}

\newcommand{\Dedications}{%
    \chapter*{Dedications}
    \addcontentsline{toc}{chapter}{Dedications} % Aggiunge 'Dedications' all'indice
}

\newcommand{\paginavuota}{
    \newpage
    \thispagestyle{empty}
}

% Formato per subsubsection
\titleformat{\subsubsection}[block]
  {\normalfont\normalsize\bfseries} % Stile del titolo
  {\thesubsubsection} % Aggiungi il numero di subsubsection
  {1em} % Spaziatura tra numero e titolo
  {} % Comando dopo il titolo

\renewcommand\thesubsubsection{\thesection.\arabic{subsection}.\arabic{subsubsection}} % Formato della numerazione per subsubsection

% Formato per paragraph
\titleformat{\paragraph}[runin]
  {\normalfont\normalsize\bfseries} % Stile del titolo
  {\theparagraph} % Aggiungi il numero di paragraph
  {1em} % Spaziatura tra numero e titolo
  {} % Comando dopo il titolo

\renewcommand\theparagraph{\thesubsubsection.\arabic{paragraph}} % Formato della numerazione per paragraph

% Definisci un nuovo comando per la subsubsubsection
% Definisci un nuovo comando per la subsubsubsection
\newcommand{\subsubsubsection}[1]{%
  \paragraph{#1}\mbox{}\\ % Usa la formattazione del paragraph e va a capo
  \addcontentsline{toc}{paragraph}{#1} % Aggiungi al sommario
}



% change this configuration with your info
% if you need fewer or more supervisors you have to change \relatore command by adding or removing lines in the table in toptesi_config
\newcommand{\thesistitle}{The role of the communicative prior on collaborative artificial agents in language games}
\newcommand{\thesisuniversitylogo}{images/polito_logo_2021_blu} % choose your logo
\newcommand{\thesiscandidatename}{Samuele Maria}
\newcommand{\thesiscandidatesurname}{Ferruggia}
\newcommand{\thesissupervisoronetitle}{Prof.}
\newcommand{\thesissupervisoronename}{Antonio}
\newcommand{\thesissupervisoronesurname}{Vetrò}
\newcommand{\thesissupervisortwotitle}{Dott.}
\newcommand{\thesissupervisortwoname}{Andrea}
\newcommand{\thesissupervisortwosurname}{Silvi}
\newcommand{\thesissupervisorthreetitle}{Prof.}
\newcommand{\thesissupervisorthreename}{Moa}
\newcommand{\thesissupervisorthreesurname}{Johansson}
\newcommand{\thesissupervisorfourtitle}{}
\newcommand{\thesissupervisorfourname}{}
\newcommand{\thesissupervisorfoursurname}{}
\newcommand{\thesisdate}{DICEMBRE 2026}
\newcommand{\thesiscourse}{COMPUTER ENGINEERING}
\newcommand{\thesisuniversity}{POLITECNICO DI TORINO}
\newcommand{\thesislevel}{MASTER} % master or bachelor
\newcommand{\thesiscandidatetext}{Candidate}
\newcommand{\thesissupervisortext}{Supervisors}

\addbibresource{bibliography.bib}

% to load the glossaries, not needed if using bib2gls
% for glossary entry
% @entry{bird,
%     name={bird},
%     description = {feathered animal},
%     see={[see also]{duck,goose}}
% }

% if this bib file does not work, try using \input{file.tex}
% where all the \newabbreviation commands have been inserted
% containing all the definitions

% Gls to capitalize first letter
% GLS for full uppercase
% for abbreviations also
% glsxtrshort for abbreviation
% similar for long, full, and capital configurations, add pl at the end for plurals
% glsentryshort, long, plural (referred to shorts) must be used when in section titles
% glslink to allow the link but use a different text (as for href)


% if you want to use also description for the abbreviations/acronyms, you should use bib2gls and define all the entries in a bib file, which is incompatible with Overleaf
\newacronym{PACE}{PACE}{Procedural Abstractions for Communicating Efficiently}
\newacronym{MDL}{MDL}{Minimum Description Length}
\newacronym{RQ}{RQ}{Research Question}
\newacronym{KL}{KL}{Kullback--Leibler}
\newacronym{EC}{EC}{Emergent Communication}
\newacronym{DSL}{DSL}{Domain-Specific Language}
\newacronym{RL}{RL}{Reinforcement Learning}
\newacronym{REINFORCE}{REINFORCE}{REward Increment = Nonnegative Factor times Offset Reinforcement times Characteristic eligibility}
\makeglossaries

\begin{document}

\overfullrule=0.00001pt % latex shows a black bar for overfulls over this dimension

%\emergencystretch=1em % to allow some stretching in the lines to avoid overfull boxes, also in bibliography, eventually can be used only before bibliography or in the preamble for the whole document, not needed if using biblatex configuration in most cases



\title{\vspace*{-5mm}\textbf{\thesistitle}\\Summary} % vspace is needed to shift upwards the title
\date{\thesisdate}
\author{\textbf{Candidate}:\\\thesiscandidatename~\thesiscandidatesurname\\
\textbf{Supervisors}:\\\thesissupervisoronetitle~\thesissupervisoronename~\thesissupervisoronesurname\\
\thesissupervisortwotitle~\thesissupervisortwoname~\thesissupervisortwosurname\\
\thesissupervisorthreetitle~\thesissupervisorthreename~\thesissupervisorthreesurname}


\maketitle

\section*{Motivation and objectives}

This thesis investigates how the distribution of communicative experience shapes the vocabularies developed by collaborative artificial agents~\cite{lazaridou_emergent_2020}. It builds on PACE (Procedural Abstractions for Communicating Efficiently)~\cite{thomas_pace_2025}, a neuro-symbolic architecture in which two neural agents learn to communicate while a symbolic mechanism introduces reusable abstractions. Unlike human participants~\cite{mccarthy_learning_2021}, these agents have no prior visual or linguistic training, so the formation of conventions can be studied in a controlled setting. The focus is on the communicative prior, understood as the frequency distribution of the scenes encountered during training.

Whereas the original PACE setup uses uniform input, this work introduces Zipfian distributions~\cite{zipf_human_1949} to vary how strongly experience concentrates on a few scenes. Three research questions guide the analysis: how concentration affects the internal structure of the emergent vocabulary; whether independently trained populations develop different vocabularies under the same concentration profile; and how these differences affect communication when agents meet unfamiliar partners.

\section*{Architecture and experimental design}

PACE implements an architect--builder game~\cite{mccarthy_learning_2021} on a $9\times9$ grid. A speaker guides a listener in reconstructing a target scene by communicating the actions of a drawing program. Initially, programs use only horizontal and vertical block-placement primitives. During training, recurring subsequences can become abstractions, allowing several primitive actions to be represented by a single reusable symbol. The architecture alternates program selection, neural communication training, and abstraction discovery. Candidate abstractions are ranked using a minimum-description-length (MDL) criterion that balances sub-sequence length against frequency, and the top-scoring candidate is added to the vocabulary. During joint training, an $\varepsilon$-greedy bandit~\cite{sutton_reinforcement_1998} decides, scene by scene, whether to use a program that contains adopted abstractions or one built entirely from primitives. The resulting symbolic libraries make the learned conventions directly inspectable.

The dataset contains 961 scenes, corresponding to all ordered pairs of 31 shapes, with 930 training scenes and 31 held-out test scenes. Training probabilities follow $p(k)\propto k^{-\alpha}$, where $k$ is a scene's position in a seed-specific ranking. At $\alpha=0$, input is uniform; increasing $\alpha$ concentrates exposure on the highest-ranked scenes. The campaign comprises 420 runs: 21 concentration values from 0.00 to 3.00, in increments of 0.15, and 20 seeds per value. Architecture and hyperparameters remain fixed. Seeds determine both neural initialisation and scene-ranking permutations, so variation across seeds combines both sources and does not isolate ranking effects.

The analysis measures abstraction adoption, program compression, held-out reconstruction accuracy, and overlap between independently learned libraries. A cross-population pipeline then pairs speakers and listeners from different runs and tracks their adaptation during joint training. Additional experiments vary speaker adaptation and exploration settings, and examine mixed input distributions and distribution shifts.

\section*{Main findings}

\paragraph{Vocabulary development is non-monotonic.}
Concentration does not simply reduce vocabulary size. Instead, three regimes emerge. Near-uniform input provides broad coverage but weak frequency preferences, so abstraction adoption is unstable across seeds. Intermediate concentration combines repeated exposure to useful patterns with enough diversity to sustain discovery, and produces the largest and most stable libraries. At high concentration, only a few scenes receive substantial exposure, which restricts the candidate pool, shrinks the libraries, and slows compression within the same training budget. Adoption also reflects more than raw frequency: the adopted abstractions are those with the highest compressive value, which combines recurrence with length. This echoes the view that lexicons balance informativeness against cost~\cite{gibson_how_2019}.

\paragraph{Shared concentration does not imply shared conventions.}
The overlap between the libraries of independently trained populations decreases as concentration increases. At low and moderate concentration, populations converge on a shared structural core beyond chance agreement; at high concentration, they develop specialised libraries tied to their own frequent scenes. These observations support a form of dialect formation: agents exposed to the same distributional shape can develop different symbolic conventions because their experience emphasises different content. Conventions diverge in a similar way between human dyads playing the same game~\cite{mccarthy_learning_2021}. The term describes population-specific libraries in this task, rather than a claim that the resulting codes reproduce human dialects.

\paragraph{Cross-population adaptation depends on the training regime.}
When partners from different populations are paired, adaptation difficulty depends mainly on the shared concentration regime rather than on how far apart the partners' training data are. Adaptation is fast at low concentration, slower at intermediate concentration, and unsuccessful within the evaluation horizon at high concentration. Pairings across different concentration values reveal an asymmetry. Speakers trained on broad input communicate more robustly with unfamiliar listeners, whereas highly concentrated speakers use vocabularies that transfer poorly. Allowing the speaker to adapt neurally does not remove this gap in the tested configurations. This points to a bottleneck in the coverage and compatibility of the learned symbolic libraries, in line with the partner-specificity of human conventions~\cite{brennan_conceptual_1996}. Distribution-shift case studies further suggest that neural recovery and symbolic reorganisation can follow different trajectories.

\section*{Conclusions and limitations}

The principal contribution is a controlled empirical framework linking input concentration to vocabulary development, cross-population divergence, and adaptation. The findings suggest that reusable conventions require both repetition and diversity, while successful communication within a population does not guarantee transfer to unfamiliar partners. PACE's explicit abstractions help distinguish changes in symbolic vocabulary from changes in neural communication.

These conclusions remain specific to the architecture, finite scene inventory, and training horizon studied. Several cross-population modifications are pilot-scale, the cross-concentration experiments use two speaker seeds, and some distribution-shift analyses use a single seed. Longer adaptation could also change outcomes classified as unsuccessful within 20 epochs. Future work should broaden seed coverage, extend training, measure vocabulary usage more fully, and separate initial intelligibility from subsequent learning. Richer environments and networks of interacting populations would test whether the same balance between shared structure and local specialisation persists beyond this setting.

\printbibliography

\end{document}
